%! Probability Models — Unit 3, Lesson 3.2 — Lesson Plan
% GRADUAL-RELEASE plan, timing 5 / 45 / 5 / 5: warm-up / guided notes and practice / spoken debrief
% / close and ASSIGN the homework (not started in class). I do / we do in rows 1-4 (problems 1-10,
% ~25 min) -> we do the Guided Practice row (problems 11-14, ~20 min). No group activity, no exit
% ticket, no independent practice set in the notes.
% Folds in the old Lessons 3.2 and 3.3 (Unit 3 condensed 2026-11-09).
\documentclass[10pt]{article}
\usepackage{saar-article}
\usepackage{saar-boxes}
\usepackage{graphicx}   % -article does NOT load graphicx
\graphicspath{{images/}}
\newcommand{\TallMath}[1]{$\displaystyle #1\rule[-1.4em]{0pt}{3.2em}$}
\newcommand{\UnitNumberName}{Unit 3: Categorical Data and Probability \quad}
\newcommand{\LessonNumberName}{Lesson 3.2: Probability Models}

\begin{document}
\begin{center}
    {\Large\bfseries \CourseName \\
    {\small\bfseries \UnitNumberName \LessonNumberName}}
\end{center}

\begin{tcolorbox}[colback=frost, colframe=cerulean, boxrule=0.9pt, arc=2mm,
    left=2mm, right=2mm, top=1.5mm, bottom=1.5mm]
    \textbf{Primary Objective:} Students will list a sample space and an event, find a theoretical
    probability and predict an expected count from it; use the complement rule and tell
    complementary from mutually exclusive events in context; organize one situation as a tree
    diagram and a two-way table; and read a Venn diagram for intersection, union, ``neither,'' and
    ``only,'' justifying why adding two circle totals overcounts ``or.''
    \\[2pt]
    \textbf{Standards:} PS.P.1a (complementary and mutually exclusive events only), PS.P.1b;
    AFDA.DA.3a, AFDA.DA.3c, AFDA.DA.3e \quad$\cdot$\quad carries forward Lesson~3.1 (two-way
    tables and relative frequency)
    \\[2pt]
    \textbf{Lesson model:} \emph{Gradual release --- I do, we do, you do.} The notes deliver the
    instruction \textbf{row by row in one Main Ideas / Notes table} (four rows, problems 1--10: you
    read the definition, mark up the display, and work problems 1, 3, 5, and~7; the class works
    the rest). The Guided Practice row (problems 11--14) is worked together with students holding
    the pen. \textbf{The homework is the individual practice}, assigned at Close \& Assign and done
    outside class. The debrief is a \emph{spoken} cold check and one board item.
    \\[2pt]
    \textbf{Scope on the page:} ``or'' is \emph{counted region by region}; the addition-rule
    formula and multiplying along a tree's branches are Lesson~3.3. The tree today carries given
    counts only. Independence and conditional probability belong to Lessons 3.3 and~3.4.
\end{tcolorbox}

\renewcommand{\arraystretch}{1.4}
\boxguard
\begin{skillbox}[Learning Targets \& Key Understandings]{goldbox}
\begin{tabularx}{\linewidth}{@{}X|X@{}}
\textbf{Learning targets (student language)} & \textbf{Key understandings (the ``why'')} \\[2pt]
\hline
\vspace{-6pt}
\begin{itemize}[leftmargin=1.1em, nosep, after=\vspace{2pt}]
  \item I can list a \textbf{sample space} and an \textbf{event}, find a theoretical probability,
        and predict an expected count. \hfill \textit{AFDA.DA.3a}
  \item I can use the \textbf{complement} rule and tell complementary from \textbf{mutually
        exclusive} events, with examples in context. \hfill \textit{AFDA.DA.3e, PS.P.1a}
  \item I can organize one situation as a tree diagram and a two-way table and read
        probabilities from each. \hfill \textit{PS.P.1b, AFDA.DA.3c}
  \item I can read a Venn diagram for intersection, union, ``neither,'' and ``only.''
        \hfill \textit{PS.P.1b}
\end{itemize}
&
\vspace{-6pt}
\begin{itemize}[leftmargin=1.1em, nosep, after=\vspace{2pt}]
  \item \textbf{Target misconception:} reading a whole circle as the ``only'' region, and counting
        the overlap twice for ``or.'' \textbf{The rule that replaces it:} a circle's total
        includes the overlap; to count ``or,'' add the regions inside the circles, each one once.
  \item \textbf{Second misconception:} counting category names instead of equally likely outcomes
        (four prize names, so $\frac14$). The denominator is the sample space, and the sectors
        are what is equally likely.
  \item Complementary events are always mutually exclusive; mutually exclusive events are
        complementary only if together they cover the whole sample space.
  \item An expected count is a prediction, not a promise.
\end{itemize}
\end{tabularx}
\end{skillbox}

\boxguard[24]
\begin{skillbox}[{Vocabulary, Concepts, \& Theorems --- taught directly in the guided notes}]{greenbox}
\small
\renewcommand{\arraystretch}{1.35}
\begin{tabularx}{\linewidth}{@{} >{\bfseries}p{3.4cm} X @{}}
Sample space & The set of all possible outcomes of a chance process. \\
Event & Any set of outcomes from the sample space. \\
Theoretical probability & With equally likely outcomes, outcomes in the event over outcomes in the sample space. \\
Expected count & $P(\text{event}) \times$ the number of trials: a prediction. \\
Complement & All the outcomes not in the event: $P(\text{not } A) = 1 - P(A)$. \\
Mutually exclusive & Two events that cannot both happen on the same outcome. \\
Complementary & Mutually exclusive \emph{and} together covering the whole sample space. \\
Intersection / union & The overlap ($A$ and $B$) / everything inside either circle ($A$ or $B$). \\
Only / neither & A circle without the overlap / outside both circles. \\
\end{tabularx}
\end{skillbox}

% THE PHASE TABLE IS A CONTRACT: 5 / 45 / 5 / 5 — I do ~25 of the 45, Guided Practice ~20.
\begin{fixedskillbox}[Lesson at a Glance --- \MeetingLength]{frost}
\renewcommand{\arraystretch}{1.25}
\begin{tabularx}{\linewidth}{@{}l c X X@{}}
\textbf{Phase} & \textbf{Min} & \textbf{Students} & \textbf{Teacher} \\
\hline
Warm-Up & 5 & Simplify fractions; fraction to decimal to percent; $\frac38$ of $200$ --- silent &
  Debrief item~3 aloud; say that these exact numbers come back in the notes \\
Guided Notes \& Practice & 45 & Fill the vocabulary box as each term is named; vote on Dani's
  count; work rows 1--4 (problems 1--10); then the Guided Practice, \emph{The Team Roster}
  (11--14), together & \textbf{I do / we do, row by row}: read the definition, mark up the display,
  work problems 1, 3, 5, 7; the class works the rest (25) $\to$ \textbf{we do} Guided Practice
  11--14 (20) \\
Debrief & 5 & Pens down --- answer aloud, nothing to fill in & Put problem~13 back on the board
  ($24$ vs.\ $20$); ask the cold check \\
Close \& Assign & 5 & Write down the homework, the due date, and item~5 as the first item &
  Announce the homework and its form (packet), preview Lesson~3.3 \\
\end{tabularx}
\end{fixedskillbox}

\begin{fixedskillbox}[Warm-Up --- Activate Prior Knowledge (5 min)]{frost}
\begin{tabularx}{\linewidth}{@{}X X@{}}
\begin{minipage}[t]{\linewidth}\vspace{0pt}
    \textbf{The three items and what each seeds}
    \begin{itemize}[leftmargin=1.1em, itemsep=1pt, topsep=2pt]
      \item \textbf{1.} Simplify $\frac{28}{40}, \frac{12}{40}, \frac{10}{40}, \frac{80}{120}$.
            \textbf{Seeds:} problem~9's $\frac{7}{10}$, problem~7's $\frac{3}{10}$, problem~8's
            $\frac14$, problem~6's $\frac23$.
      \item \textbf{2.} $\frac38 = 0.375 = 37.5\%$; $\frac58 = 0.625 = 62.5\%$; $\frac{7}{10} =
            0.7 = 70\%$. \textbf{Seeds:} problem~1's three forms and problem~3's complement.
      \item \textbf{3.} $\frac38$ of $200$ is $75$; $1 - \frac38 = \frac58$. \textbf{Seeds:}
            problem~2's expected count and the complement rule in row~2.
    \end{itemize}
\end{minipage}
&
\begin{minipage}[t]{\linewidth}\vspace{0pt}
    \textbf{Running it}
    \begin{itemize}[leftmargin=1.1em, itemsep=1pt, topsep=2pt]
      \item \textbf{Debrief item~3 aloud.} ``Three eighths of something'' and ``what is left of
            one whole'' are the two moves of the first two rows.
      \item Items 1--2 are arithmetic only: do not let them expand. Read the answers once.
      \item Note who struggles to simplify; those students can leave a probability unsimplified
            in the notes and still follow the idea.
    \end{itemize}
\end{minipage}
\end{tabularx}
\end{fixedskillbox}

\boxguard
\begin{skillbox}[Hook --- before anyone sits down]{frost}
The hook is on \textbf{page 1 of the notes}, under the vocabulary box, and on the board:
\textit{``$40$ students were surveyed. $18$ bought a taco and $15$ bought a lemonade. Dani adds:
$18 + 15 = 33$ bought a taco or a lemonade.''} Students circle \emph{right}, \emph{too high},
\emph{too low}, or \emph{can't tell yet} and write one line of reason \emph{before} anyone speaks;
then take the hand vote and write the split on the board. Most rooms vote \emph{right}.
\textbf{Do not settle it.} Say only that the answer is problem~9, in the \emph{Venn Diagram} row,
and that they will vote again there.
\end{skillbox}

\boxguard
\begin{skillbox}[Guided Notes \& Practice --- I do, then we do (45 min)]{frost}
\textbf{This is the whole lesson.} The notes are \textbf{one Main Ideas / Notes table}: four
instruction rows (problems 1--10), then the Guided Practice row (11--14). \textbf{The rhythm in
every row is the same}: read the definition aloud, mark up the display, work the first problem
yourself, then hand the rest to the room. The vocabulary box is filled \emph{as each term is
named}. The only blanks are the eight cells of the row~3 table.
\begin{multicols}{2}
\textbf{I do / we do, row by row (about 25 min).}

\emph{Row 1, Theoretical Probability (problems 1--2).} Name \emph{sample space} and \emph{event}
on the prize wheel. Work problem~1 yourself: Kira's $\frac14$ counts four prize names, but the
eight sectors are what is equally likely, so $P(D) = \frac38 = 0.375 = 37.5\%$. The class does 2:
$75$ free drinks of $200$ spins, ``about'' because the spins are random (AFDA.DA.3a). Five minutes.

\emph{Row 2, Complement Rule (problems 3--4).} Two wheels side by side, students label them.
Work problem~3 yourself ($\frac58$, then $\frac34$). The class does 4: free drink and free dessert
are mutually exclusive only (sectors 4, 5, 8 are left over); free drink and even-numbered sector
are neither (sector~6 is both). Ask aloud for one festival example of each kind (AFDA.DA.3e,
PS.P.1a). Seven minutes.

\emph{Row 3, Tree and Two-Way Table (problems 5--6).} The same $120$ orders twice; students copy
the branch ends into the table. Work problem~5 yourself (burger, then water: $\frac{15}{120} =
\frac18$). The class does 6 ($\frac{80}{120} = \frac23$). \textbf{Read counts off the tree; do not
multiply along branches}, which is Lesson~3.3 (PS.P.1b, AFDA.DA.3c). Five minutes.

\emph{Row 4, Venn Diagram (problems 7--10) --- the crux.} Name the four regions on the survey and
work problem~7 yourself ($\frac18$ and $\frac{3}{10}$). \textbf{Problem~8 is the trap}: Reese reads
the whole circle, $15$, as ``only lemonade''; the region is $10$. Then \textbf{problem~9 is the
crux: re-take the hook vote before anyone counts.} Region by region, $13 + 5 + 10 = 28$, and
$P = \frac{28}{40} = \frac{7}{10}$, not Dani's $33$. Problem~10 names the $5$ counted twice and
checks with $1 - \frac{12}{40}$. Land the boxed sentence word by word. Eight minutes.

\columnbreak
\textbf{We do (about 20 min).} The Guided Practice row, \emph{The Team Roster} (problems 11--14),
a fresh context where the diagram prints only the circle totals. \textbf{Students hold the pen; you
hold the questions.} Problem~11 fills the regions ($11$, $4$, $5$, $10$). Problem~12 reads
$\frac16$ and $\frac13$. \textbf{Problem~13 is the crux transferred}: take a hand count on Coach's
$24$ before anyone counts; the union is $20$, $P = \frac23$. Problem~14 ties back to row~2: not
mutually exclusive ($4$ play both); ``soccer or basketball'' and ``neither'' are complementary
($20 + 10 = 30$).

Circulate during Guided Practice and demand the reason:
\begin{itemize}[leftmargin=1.1em, nosep]
  \item \textbf{Problem 11:} ``Does the $15$ include the $4$? So what goes in the left region?''
  \item \textbf{Problem 13:} ``Point to every athlete you counted. Did your finger land on the
        overlap twice?''
  \item \textbf{Problem 14:} ``What would the diagram look like if the two sports \emph{were}
        mutually exclusive?'' (Circles that do not touch.)
\end{itemize}
While circulating, pick the one paper to put up in the debrief --- a problem~13 that counted
region by region and named the $4$. \textbf{This circulation is your formative read}.
\end{multicols}
\end{skillbox}

\boxguard[24]
\begin{skillbox}[Debrief --- whole class, spoken (5 min)]{frost}
Pens down, papers up. \textbf{This debrief is spoken.}
\begin{multicols}{2}
\begin{enumerate}[leftmargin=1.3em, nosep, itemsep=3pt]
  \item \textbf{The one board item: problem~13.} Put up the paper you picked ($11 + 4 + 5 = 20$,
        not $24$) and make a student say the sentence: \emph{adding the circle totals counts the
        overlap twice; to count ``or,'' add each region once.}
\end{enumerate}
\columnbreak
\textbf{Check for understanding --- ask it cold, whole class.} \emph{``$32$ students: $16$ are in
band, $10$ in choir, and $6$ in both. How many are in band or choir, and how many in only choir?''}
Thirty seconds of think time, then cold-call three students. Correct answer: \textbf{$20$ in band
or choir} ($10 + 6 + 4$), not $26$; \textbf{$4$ in only choir}, not $10$. \emph{This replaces the
exit ticket.}

\textbf{Formative read.} Sort the answers into three piles: \textbf{(a)} $20$ and $4$;
\textbf{(b)} one of the two right; \textbf{(c)} $26$ and $10$ --- both halves of the
misconception. \textbf{If pile (c) runs past a third of the room, open Lesson~3.3 by re-running
problems 8--9}, and say so plainly.
\end{multicols}
\end{skillbox}

\boxguard
\begin{skillbox}[Homework --- scored 1 to 5, assigned today, due after two study halls]{goldbox}
Six items on \textbf{a bookstore loyalty program}: \textbf{1} a rewards wheel, $P(\$3\text{ off})
= \frac{3}{10}$, and the \emph{interpret} item, the meaning of about $150$ of $500$ spins;
\textbf{2} the complement ($\frac{7}{10}$) and classifying \$1 off and \$3 off (mutually exclusive
only); \textbf{3} and \textbf{4} are the \textbf{contrast pair} on one Venn diagram: bought a
mystery ($26$, $\frac{13}{30}$) against bought \emph{only} a mystery ($19$, $\frac{19}{60}$);
\textbf{5} \emph{the crux head-on and the justification}: Priya's $26 + 17 = 43$, the true $36$,
and who was counted twice; \textbf{6} \emph{the spiral item}, Lesson~3.1's relative frequency in
a two-way table ($0.55$ and, among teens, $0.67$).

\textbf{Assigned at Close \& Assign and done outside class. Scoring:} 12 points, 2 per item, which
the grading report turns into the 1-to-5 score. \textbf{Form: packet.} The packet pages are this
lesson's form; every final answer goes in its box. The teacher may switch the form at Close \&
Assign.

\textbf{Preview of the next class.} Lesson~3.3, Addition and Multiplication Rules: today's
region-by-region count becomes a rule, and the tree gets probabilities on its branches.

\textbf{Connections \& Big Ideas.} \textit{Unit 3 core idea:} a probability model organizes the
outcomes so that counting them is honest. \textit{Standards covered:} PS.P.1a (complementary and
mutually exclusive events; dependent and independent events wait for 3.3--3.4), PS.P.1b (Venn
diagrams, tree diagrams, two-way tables), AFDA.DA.3a (theoretical probability and predictions),
AFDA.DA.3c (representing and calculating with these models), AFDA.DA.3e (defining and giving
contextual examples). \textit{Reaches back to:} Lesson~3.1's two-way tables. \textit{Spirals
forward to:} Lesson~3.3's addition and multiplication rules and Lesson~3.4's conditional
probability.
\end{skillbox}

\boxguard
\begin{skillbox}[Watch For (while circulating)]{redbox}
\begin{itemize}[leftmargin=1.2em, nosep]
  \item \textbf{Whole circle read as ``only''} (problems 8 and~11, HW~3--4 --- the target
        misconception). Probe: \emph{``Does that total include the people in the overlap?''}
  \item \textbf{Overlap counted twice for ``or''} (the hook, problems 9 and~13, HW~5 --- the
        target misconception). Probe: \emph{``Touch each region once. What total did you get?''}
  \item \textbf{Counting names, not outcomes} (problem~1). Probe: \emph{``Are those four things
        equally likely? What on the wheel is?''}
  \item \textbf{Complementary and mutually exclusive treated as the same} (problems 4 and~14,
        HW~2). Probe: \emph{``Is anything left over?''}
  \item \textbf{An expected count read as a guarantee} (problem~2, HW~1). Probe: \emph{``Would
        you bet on exactly $75$?''}
  \item \textbf{Multiplying branch numbers} (problem~5). Probe: \emph{``The path already ends in
        a count. What does it say?''}
\end{itemize}
Cold-call prompts: \emph{``Name two events at this school that cannot both happen.''} \quad
\emph{``When \emph{is} adding the two totals correct?''} (When the circles do not overlap.)
\end{skillbox}

\boxguard
\begin{skillbox}[Close \& Assign (5 min)]{goldbox}
\textbf{The homework is assigned here and done outside class --- it is not started in class.}
Hand the launch: six items on a bookstore loyalty program, \textbf{packet form}, scored 1 to 5, due
at the start of the first class after two study halls. The teacher may switch the form here.
Point to the \emph{Keep in Mind} box and rows 2 and~4 of the notes. \textbf{Tell students to do
item~5 first if they are short on time}, and that items 3 and~4 are a pair: read both before
answering either. Then close on what changed: \emph{``You walked in with fractions. You walk out
with a sample space, a complement, and a Venn diagram you can count without counting anyone
twice.''} \textbf{Preview:} Lesson~3.3, Addition and Multiplication Rules.
\end{skillbox}

% ── Teacher notes — one per component, in packet order (three) ──────────────
\begin{teachernote}[Warm-Up]
Five minutes, silent, timed. Every number here returns in the notes, so say that rather than
reviewing at length. Item~3 is the one to debrief: $\frac38$ of $200$ becomes the expected count
in problem~2, and $1 - \frac38$ becomes the complement rule in problem~3. A student who cannot
simplify can still follow the lesson with unsimplified fractions; do not stop for it.
\end{teachernote}

\begin{teachernote}[Guided Notes \& Practice]
\textbf{Pace it 25 / 20, then 5 for the debrief and 5 to assign the homework.} You work 1, 3, 5,
and~7; the class works the rest. Rows 1 and~3 must move, five minutes each.
\textbf{Spend the time in row~4.} The trap (problem~8) comes before the crux (problem~9) on
purpose: a student who has just separated a circle total from its ``only'' region can see why the
circle totals overlap. \textbf{Re-take the hook vote at problem~9 before anyone counts}; if
students count first, the vote is empty. The boxed sentence comes after 9 and~10, so the rule
arrives as the answer to the vote. \textbf{Do not state the addition rule}: students who offer
``subtract the overlap'' are right, and Lesson~3.3 makes it a rule; today every ``or'' is counted
region by region. On the tree, read counts only.

\textbf{Guided Practice (11--14) is where the pen changes hands.} The diagram prints only circle
totals, so the student has to separate each ``only'' region before anything else. If time is short,
protect 11 and~13; problem~12 can be said aloud. If the misconception pile is more than a third of
the room, \textbf{open Lesson~3.3 by re-running problems 8--9}.
\end{teachernote}

\begin{teachernote}[Homework]
\textbf{Scored 1 to 5 from 12 points (2 per item), due at the start of the first class after two
study halls; assigned at the close, not started in class.} Every part is boxed: 1a, 2a, 2b, 3a,
3b, 4a, 4b, 5a, 6a, and~6b are short boxes worth 1 point each; 1b and~5b are written responses
worth 1 point each, scored by their rubrics. \textbf{Item~5 is the one to read first.} Sort it:
named the $7$ counted twice (ready); got $36$ with no reason (needs the language); wrote $43$ (the
misconception survived). Items 3 and~4 are a diagnostic pair: the same answer on both means the
student is reading the circle as the region. If more than a third land in the last pile, open
Lesson~3.3 with problems 8--9. \textbf{Packet is the form for this lesson}; the teacher may switch
it at Close \& Assign.
\end{teachernote}
\end{document}
